\section{圈算子与线算子}

\label{sec:loopop} 

\begin{keytakeaways}
\begin{itemize}[leftmargin=1.4em,itemsep=2pt,topsep=2pt]
\item \textbf{物理图像}：本节讨论的四维 TQFT 允许插入与定向一维子流形 ${\cal S}\subset M$ 相联系的算子。这些算子的参数不是普通的「点」，而是环路或通向无穷远的线，因而被称为圈算子与线算子。
\item \textbf{数学工具}：这类算子依赖 ${\cal S}$ 的同伦类而不只是同调类；在 $t=\pm i$ 处，可以用 $A$ 与标量场 $\phi$ 的复化联络 $\CA=A+i\phi$（或 $\overline\CA=A-i\phi$）的 holonomy 构造具有拓扑对称性的 Wilson 算子。
\item \textbf{关键结论}：Wilson 算子是四维理论中最基本的圈算子；它的磁对偶 't Hooft 算子将在第~\ref{sec:tophooft}~节引入。紧化到二维后，线算子可以被解释为作用在膜上的算子，这正是通向几何 Langlands 纲领的入口。
\item \textbf{边界与展望}：对有端点或通向无穷远的线算子须小心处理；在非紧四维流形上，静态线算子必须纳入量子化以构造物理态空间，这是本文应用中最重要的情况。
\end{itemize}
\end{keytakeaways}

与第~\ref{sec:twisting}~节引入的四维 TQFT 表面上来源相似的其他理论相比，这些 TQFT 真正特别之处在于：它们允许与定向一维流形 ${\cal S}\subset M$ 相联系的算子，而且这些算子依赖 ${\cal S}$ 的同伦类，而不只是同调类。要点并不是我们想研究 $M$ 的基本群，而是这些一维流形算子在关于膜的问题上被证明具有非常有趣的性质。

若 $M$ 是紧的且没有边界，则要求 ${\cal S}$ 是闭圈。为反映这一情形，一维流形算子通常被称为圈算子。然而更一般地，在非紧四维流形上，我们允许通向无穷远的线算子。一个重要的例子是静态线算子，其轨迹在空间中某个固定点处贯穿全部时间。在量子化理论以构造物理态空间时，必须把这样的线算子包括进来。这实际上将是我们应用中最重要的情况。\footnote{因此，其他情形——例如终止在膜上的线算子——虽然在一般情况下有意义，但在本文中并不重要。}

% 说明：以下四行是脚手架把原文的 \nref 文献定义转成的 \cite；原文此处并不产生引用标记，
% 这些文献在下文已有正常引用。为保持 \cite 计数与 skel 一致，此处按注释保留。
% \cite{ref84}
% \cite{ref85}
% 
% 

在第~\ref{sec:topwil}~节，我们讨论最基本的圈算子或线算子，通常称为 Wilson 算子。在与本文所考虑的理论类似的 TQFT 中，Wilson 算子最早由 Blau 与 Thompson 构造 ，其中的联络被一个标量场修改以保证拓扑对称性。类似的规范不变算子也被用于研究五维超对称杨--米尔斯理论 。这类算子在弦论的当代进展中是重要的（连同它们的磁对偶，我们将在第~\ref{sec:tophooft}~节引入）。它们最早是在半 BPS 情形下被引入的 \cite{ref86,ref87}。我们所考虑的拓扑线算子，与 \cite{ref88} 中对任意环路 ${\cal S}\subset \Bbb{R}^4$ 定义的 1/16 BPS 线算子最为相似。

\begin{yizhu}[含糊点拨]
按：「依赖同伦类而非同调类」是本节标题所示算子的关键特征：算子的插入位置是 $M$ 中带定向的一维子流形，拓扑对称性使关联函数只依赖其同伦类。原文特意声明这并不是要研究 $\pi_1(M)$——研究基本群本身没有意义，有意义的是这类算子与膜的关系，后文会看到它们变成膜范畴上的函子。
\end{yizhu}

在第~\ref{sec:tophooft}~节，我们描述对偶的 't Hooft 算子。在第~\ref{sec:comptwo}~节，我们讨论这类算子在紧化到二维后的解释。最后，在第~\ref{sec:lineop}~节，我们把线算子作为膜上的算子来研究。

\subsection{拓扑 Wilson 算子}

\label{sec:topwil}

在规范场论中，最基本的圈算子是 Wilson 圈算子。我们考虑规范群为 $G$、$G$-丛 $E\to M$ 上的联络为 $A$ 的规范场论。令 ${\cal S}$ 是 $M$ 中的定向闭圈，其定向比如说由选取一个一形式 $ds$ 给出，其中 $s$ 是沿 ${\cal S}$ 的参数。我们取定 $G$ 的一个不可约表示 $R$，并令\footnote{由于我们的协变导数为 $D=d+A$，holonomy 算子在指数上带一个负号。用 hermitian 规范场 $A'=iA$（回忆式 \eqref{eq:boloko}）表述时，holonomy 算子为 ${\rm Tr}_R\,P\,\exp\left(i\oint_{\cal S} A'\right)$。}
% [WARN] 公式 \eqn\eqref{eq:ytro} 无编号或重复
\begin{equation*}
\tag*{(6.1)}
W_0(R,{{\cal S}})={\rm Tr}_R
\,P\,\exp\left(-\oint_{{\cal S}} A\right).
\end{equation*}
换言之，$W_0(R,{\cal S})$ 是 $A$ 绕 ${\cal S}$ 的 holonomy 在表示 $R$ 中的迹。

\begin{yizhu}[易错点]
按：式 (6.1) 指数上的负号来自协变导数约定 $D=d+A$（见脚注），换成 hermitian 规范场 $A'=iA$ 后就是通常的 ${\rm Tr}_R P\exp(i\oint A')$。这里的 $W_0$ 还是普通（未扭化）的 Wilson 圈算子；它一般不具有拓扑对称性，正是下面要修改的对象。
\end{yizhu}

现在我们回看第~\ref{sec:twisting}~节中的式 \eqref{eq:bono}，其中我们确定了 $A$ 在拓扑对称性下的变换规律：% 式 (6.2)
\setcounter{equation}{1}
\begin{equation}\label{eq:likp}
\delta_T
A_\mu=iu\psi_\mu+iv\tilde\psi_\mu.
\end{equation}
显然，对任意非零的 $(u,v)$，Wilson 圈算子不是不变的。然而，扭化的 ${\EUN}=4$ 杨--米尔斯理论还有一个伴随值的一形式 $\phi=\sum_{\mu=0}^3\phi_\mu\,dx^\mu$，其变换为
% 式 (6.3)
\setcounter{equation}{2}
\begin{equation}\label{eq:nikp}
\delta_T\phi_\mu=iv\psi_\mu-iu\tilde\psi_\mu.
\end{equation}
可以通过给联络 $A$ 加上 $\phi$ 的一个倍数来修改它。

确切地说，当 $t=v/u$ 等于 $\pm i$ 时，$A$ 与 $\phi$ 的一个线性组合是不变的。对 $t=i$，不变的组合是 $\CA=A+i\phi$；在 $t=-i$ 时，它是 $\overline\CA=A-i\phi$。于是我们把超对称 Wilson 圈算子定义为 $\CA$ 或 $\overline\CA$ 绕 ${\cal S}$ 的 holonomy。因此
% 式 (6.4)
\setcounter{equation}{3}
\begin{equation}\label{eq:ymop}
W(R,{\cal S})=\Tr_R\,P\exp\left(-\oint_{{\cal
S}}\CA\right)=\Tr_R\,P\exp\left(-\oint_{{\cal
S}}(A+i\phi)\right),~~ t=i
\end{equation}
以及 % 式 (6.5)
\setcounter{equation}{4}
\begin{equation}\label{eq:jymop}
W(R,{{\cal S}})=
\Tr_R\,P\exp\left(-\oint_{{\cal S}}{\overline\CA}\right)
=\Tr_R\,P\exp\left(-\oint_{{\cal S}}(A-i\phi)\right), ~~t=-i.
\end{equation}
$SU(4)_{{\cal R}}$ 的中心生成元 $\EUI$ 通过 $t\to -t$ 与 $\phi\to -\phi$ 作用，把其中一个公式变成另一个。

\begin{yizhu}[物理图像]
按：$t=\pm i$ 是 GL 扭化的两个特殊点：此时拓扑超荷的组合恰好使 $A\pm i\phi$ 不变，于是可以取它的 holonomy 作为拓扑不变的圈算子。$\CA=A+i\phi$ 正是 Hitchin 方程中出现的复化联络；在 $t=i$ 时拓扑方程 \eqref{eq:loopme} 断言 $\CA$ 平坦，这与后面把 Wilson 算子解释为膜上的算子密切相关。
\end{yizhu}

在 $t=\pm i$ 时，具有拓扑对称性的 Wilson 圈算子的存在是极其自然的。在这些 $t$ 值下，拓扑方程 \eqref{eq:loopme}\ 断言 $\CA$ 的平坦性。所以 $\CA$ 或 $\bar\CA$ 绕闭圈的 holonomy 是自然的不变量。同样地，若 $t\not=\pm i$，这类 holonomy 不是自然的不变量，因为拓扑方程不一定蕴含平坦性。（第~\ref{sec:vanishing}~节的消灭定理把超对称位形与一般 $t$ 下的复平坦联络联系起来，但它们需要某些整体输入，并不在所有情形下成立，例如涉及可能应用该 TQFT 的膜时。）

\begin{yizhu}[易错点]
按：$-i$ 与 $+i$ 给出的两个公式由 $\EUI$ 互换，注意不要漏掉 $t\not=\pm i$ 的例外：此时 $\CA$ 或 $\bar\CA$ 的 holonomy 不是拓扑不变量。括号内的提醒也很重要——消灭定理（把超对称位形等同于复平坦联络）依赖整体性假设，在有膜的情形可能失效，这一点在后文把线算子变成膜上算子时会反复用到。
\end{yizhu}

虽然我们是对闭圈 ${\cal S}$ 表述这些定义的，但它们在开曲线情形有很好的类比。若 ${\cal S}$ 的端点为 $p$ 与 $q$，我们就把 $W(R,{{\cal S}})$ 不定义为迹，而定义为从 $p$ 处 $E$ 的纤维（在表示 $R$ 中取）到 $q$ 处纤维的平行移动（关于联络 $\CA$ 或 $\bar\CA$）矩阵。于是 $W(R,{{\cal S}})$ 是从 $p$ 处的向量空间到 $q$ 处的向量空间的映射。它随后被纳入一个更大的量子构造，涉及初始与最终的量子态，正如我们在引言中所提到的。当我们真正进入几何 Langlands 纲领时，Wilson 算子的这个推广版本将是必不可少的。

\begin{yizhu}[物理图像]
按：开曲线情形把 Wilson 算子从「一个数」（迹）升级为「一个矩阵」（平行移动），从而可以作用在边界态空间上。这正是「线算子作为膜上的算子」的技术雏形：线算子的两个端点可以贴在膜上（或对应初始/末态），于是它的作用变成膜上态空间之间的映射，后文将把它实现为 Hecke 算子之类的函子。
\end{yizhu}

上下文清楚时，我们把 $W(R,{{\cal S}})$ 简写为 $W(R)$ 或干脆写 $W$。
